% !TEX root = ../my-thesis.tex
%
\IfLanguageName{english}{\pdfbookmark[0]{Abstract}{Abstract}}{\pdfbookmark[0]{Zusammenfassung}{Zusammenfassung}}
\IfLanguageName{english}{\addchap*{Abstract}}{\addchap*{Zusammenfassung}}
\label{sec:abstract}

This thesis presents an end-to-end pipeline that generates, structures, and
formally evaluates Abstract Argumentation Frameworks from natural-language debate
topics using large language models. Rather than treating attack detection as a
binary decision, the system prompts a language model to produce typed
attack candidates, each annotated with an attack type, namely rebuttal,
undercutting, or undermining, a confidence score, and a natural-language
justification. A two-stage confidence filter routes candidates to acceptance,
rejection, or verification. Accepted attacks are reduced to binary facts and
handed to a Prolog solver that computes grounded, preferred, and stable
extensions under Dung semantics. Evaluated on eight topics from the UKP
Sentential Argument Mining Corpus, five of eight frameworks
admit several preferred and stable extensions, each corresponding to an
internally coherent position that a standalone language model leaves implicit.
An ablation on identical atomized argument sets shows that typed and binary
prompting extract substantially different attack relations, and that the
multiplicity of extensions arises entirely through the symmetrization of
rebuttals, which only typed extraction makes possible. An evaluation against a
manual reference annotation of 201 argument pairs on one topic, created under
written guidelines on a frozen argument set, shows that the typed extraction is
measurably closer to this annotation than the binary baseline. It
recovers three quarters of the annotated attacks and reaches a directed F1 of
0.615 against 0.458.

\IfLanguageName{english}{
\vspace*{20mm}

{\usekomafont{chapter}Abstract (German)}
\label{sec:abstract-diff}
% Set german quote style and language system.
\setquotestyle{german}
\selectlanguage{ngerman}

Diese Arbeit stellt eine durchgängige Pipeline vor, die mithilfe großer
Sprachmodelle aus natürlichsprachlichen Debattenthemen abstrakte
Argumentationsframeworks erzeugt, strukturiert und formal auswertet. Statt die
Angriffserkennung als binäre Entscheidung zu behandeln, erzeugt das System
typisierte Angriffskandidaten, die jeweils mit einem Angriffstyp (Rebuttal,
Undercutting oder Undermining), einem Konfidenzwert und einer Begründung
versehen sind. Ein zweistufiger Konfidenzfilter leitet die Kandidaten zur
Annahme, Ablehnung oder Verifikation. Angenommene Angriffe werden auf binäre
Fakten reduziert und an einen Prolog-Solver übergeben, der grounded, preferred
und stable Extensions unter Dung-Semantik berechnet. In der Auswertung auf acht
Themen des UKP-Korpus besitzen fünf von acht Themen
mehrere preferred und stable Extensions, die jeweils einer in sich
konsistenten Position entsprechen, welche ein einzelnes Sprachmodell nur
implizit lässt. Eine Ablation auf identischen atomisierten
Argumentmengen zeigt, dass typisiertes und binäres Prompting deutlich
verschiedene Angriffsrelationen extrahieren und dass die Mehrfach-Extensions
vollständig durch die Symmetrisierung der Rebuttals entstehen, die nur die
typisierte Extraktion ermöglicht. Eine Auswertung gegen eine manuell
annotierte Referenz aus 201 Argumentpaaren auf einem Thema, erstellt nach
schriftlichen Richtlinien auf einer eingefrorenen Argumentmenge, zeigt, dass
die typisierte Extraktion dieser Annotation
messbar näher kommt als die binäre Baseline. Sie findet drei Viertel der
annotierten Angriffe und erreicht einen gerichteten F1 von 0,615 gegenüber
0,458 für die Baseline.

% reset to english
\selectlanguage{english}
\setquotestyle{english}

}
{}
